\section{Introduction}

Pardon's universal curve-counting algebra expresses the enumeration
of effective cycles on a projective Calabi--Yau threefold in terms of
local curves \cite{Pardon}. Its rational generators $x_{g,m,0}$ are
defined by the equivariant theory of degree-$m$ cycles in a rank-two
bundle over a genus-$g$ curve, with scalar multiplication on the fibers.
Write $x_{g,m}=x_{g,m,0}$ and
\[
 f_g(q)=1+\sum_{m\ge1}x_{g,m}q^m.
\]
The coproduct and cycle-division operations imply a unique rational
factorization of the universal cycle series,
\begin{equation}\label{eq:factor-intro}
 \Z_X=\prod_{\beta\ne0,\,g\ge0}f_g(Q^\beta)^{e_{\beta,g}(X)},
 \qquad e_{\beta,g}(X)\in\QQ,
\end{equation}
when $H_2(X;\ZZ)$ is torsion-free. Each class has only finitely many
nonzero exponents. Conjecture~1.4 of \cite{Pardon} asserts that the
exponents in this factorization are integers. We give a counterexample.

The compact threefold used below is a general anticanonical
hypersurface in the smooth projective toric fourfold specified in
Section~\ref{sec:geometry}. It contains a surface $F\simeq\FF_1$.
Let $s$ be its negative section and $f$ a ruling fiber.

\begin{theorem}\label{thm:main}
The threefold $Y$ of Section~\ref{sec:geometry} is a smooth projective
Calabi--Yau threefold with torsion-free $H_2(Y;\ZZ)$. The class
$\alpha=3s+4f$ is primitive, and its scalar-local product exponents satisfy
\[
 e_{\alpha,3}(Y)=-14,\qquad
 e_{3\alpha,10}(Y)\in\ZZ+\frac23.
\]
\end{theorem}

The arithmetic mechanism is a general comparison formula. Denote by
$c_{\beta,g}(X)$ the integer elementary exponents in the expansion of
Ionel--Parker \cite[Theorem 8.1]{IP}. Their elementary local series
uses the antidiagonal action on a balanced local curve
\cite[Proposition 3.3]{IP}. These integers differ from the scalar
exponents in \eqref{eq:factor-intro} and, in general, from the usual
Gopakumar--Vafa numbers. We prove in Proposition~\ref{prop:congruence}
that every primitive class $\alpha$ satisfies
\begin{equation}\label{eq:obstruction-intro}
 e_{3\alpha,10}(X)
 =c_{3\alpha,10}(X)-\frac{4720}{3}c_{\alpha,3}(X).
\end{equation}
Thus integrality of the exponent on the left requires
$3\mid c_{\alpha,3}(X)$. For our class, $c_{\alpha,3}(Y)=-14$.
The unknown integer $c_{3\alpha,10}(Y)$ need not be calculated.

The proof has two parts. First, the Frobenius algebra of
Bryan--Pandharipande \cite{BP} determines the scalar local series in
degree three. We give its multiplication matrix, pairing, cap and
characteristic polynomial, retaining the signs and phases distinguishing
scalar from antidiagonal residues. The connected comparison has a
$4720/3$ coefficient in the genus-ten term. A leading-order estimate
excludes contributions from all higher local genera.

Second, nef divisors on the compact hypersurface force every cycle
in a multiple of $\alpha$, and every effective summand of $\alpha$,
onto $F$. Cycle support alone does not exclude transverse nonreduced
curve schemes. A filtration by the powers of the ideal of $F$ proves
that a stable pair in class $\alpha$ has Euler characteristic at least
$-2$. At equality its curve is a divisor in $|3s+4f|\simeq\PP^{13}$.
We identify this minimal moduli space scheme-theoretically, including
its nonreduced divisors, and obtain virtual count $-14$. The same
filtration shows that lower products in the logarithm have nonnegative
Euler characteristic. The primitive genus-three coefficient is
therefore exactly $-14$.

The geometry also occurs in the conifold specialization calculation of
\cite{Transition}. Here the needed fan, nef divisors, contracted locus
and surface identifications are supplied directly. No transition
formula is used to prove Theorem~\ref{thm:main}. The universal algebra
and local TQFT are imported from \cite{Pardon,BP}; the integer elementary
expansion is imported from \cite{IP}. The result concerns the proposed
integer product expansion in the specified scalar generators. Pardon's
integral geometric-generator description \cite[Theorem 1.2]{Pardon}
and ordinary numerical Gopakumar--Vafa integrality are compatible with
the counterexample.
